% ==========================================
% BÖLÜM 5: GELİŞTİRİCİ VE KATKI SAĞLAMA REHBERİ
% ==========================================
\chapter{Geliştirici Rehberi}

Bu bölüm, uygulamanın gelecekte yeni özelliklerle genişletilmesi, değişen kurum içi rapor yapılarına kolayca adapte edilmesi, veritabanı yapılarının güncellenmesi ve mimariye yeni dinamik panellerin eklenmesi süreçlerinde takip edilmesi gereken standartları detaylandırmaktadır.

Uygulamanın sürdürülebilir kalması için geliştiricilerin aşağıda yer alan 4 temel senaryoyu ve katmanlı mimari kurallarını adım adım uygulaması büyük önem taşımaktadır.

\vspace{0.4cm}

% ==========================================
% SENARYO 1: RAPOR BİÇİMİ / SÜTUN DEĞİŞİKLİĞİ
% ==========================================
\section{Senaryo 1: Gelen Rapor Biçimi veya Sütun Yapısı Değişirse}

Sistemdeki en büyük mimari avantaj, kod satırlarına sert biçimde yazılmış (\textit{Hardcoded}) sayfa/sütun isimlerinin veya dinamik metinlerin bulunmamasıdır. Kurum içi raporlarda sütun harfleri, sayfa isimleri veya vaka durum başlıkları değiştiğinde izlenecek adımlar aşağıda detaylandırılmıştır:

\vspace{0.2cm}

\textbf{\textcolor{primary}{1. Yapılandırma Modülüne (\texttt{mod\_Config}) Erişin:}}
\begin{itemize}\setlength{\itemsep}{2pt}
    \item Klavyeden \textbf{\texttt{Alt + F11}} kısayol tuşlarına basarak VBA Editörüne girin ve sol menüden \texttt{mod\_Config} modülünü açın.
    \item Tüm sistem sabitlerinin kategorize edildiği bu alanda sadece ilgili sabit değerini güncelleyin.
\end{itemize}

\vspace{0.15cm}

\textbf{\textcolor{primary}{2. Sütun Harflerini ve Hücre Adreslerini Revize Edin:}}
Rapor sayfalarındaki sütun yerleri değiştiğinde \texttt{mod\_Config} içerisindeki ilgili harfleri güncelleyin:
\begin{lstlisting}[caption={Sütun Sabitlerinin Güncellenmesi}]
' Eski: Public Const ALL_CASES_STATUS_COL As String = "AF"
' Yeni Sütun Konumu:
Public Const ALL_CASES_STATUS_COL As String = "AG"
Public Const ALL_CASES_CLOSE_DATE_COL As String = "X"
\end{lstlisting}

\vspace{0.15cm}

\textbf{\textcolor{primary}{3. Durum Metinlerini ve Kelime Filtrelerini Genişletin:}}
Sisteme yeni bir vaka durumu eklendiyse veya durum kelimeleri değiştiyse (Örn: "Çözüldü" yerine "Tamamlandı" veya "İşlemde" yerine "Devam Ediyor" geldiyse); sabit listelerine yeni kelimeyi virgülle ekleyin:
\begin{lstlisting}[caption={Geçerli Durum Sabitlerinin Güncellenmesi}]
Public Const VALID_CLOSED_STATUSES As String = "Çözüldü,Kapandı,Kapalı,Cozuldu,KAPANDI,KAPALI,Tamamlandı,TAMAMLANDI"
Public Const VALID_OPEN_STATUSES   As String = "Açık,Atandı,Beklemede,İşlemde,Acik,Atandi,Islemde,Devam Ediyor"
\end{lstlisting}

\vspace{0.15cm}

\textbf{\textcolor{primary}{4. İş Mantığı Katmanını Koruyun:}}
Sütun veya durum ismi değişikliklerinde \texttt{mod\_Services} içerisindeki algoritmik kodlara dokunmanıza kesinlikle gerek yoktur; arka plan servisleri dinamik olarak yenilenen bu sabitleri okuyarak kesişim hücrelerini hesaplayacaktır.

\newpage

% ==========================================
% SENARYO 2: ARAYÜZE YENİ PANEL EKLEME
% ==========================================
\section{Senaryo 2: Arayüze Yeni Bir Panel Ekleme (\texttt{CreateYEPYENIPanel})}

Login paneline veya genel iş akışına yepyeni bir işlem ekranı (Panel) dahil edilmek istendiğinde projenin Katmanlı Mimarisi sırasıyla takip edilerek adımlar yürütülür:

\vspace{0.2cm}

\textbf{\textcolor{primary}{1. Adım: Sabit ve Görsel Metin Tanımlamaları (\texttt{mod\_Config})}}
\texttt{mod\_Config} modülünde yeni eklenen panelin başlığı ve buton yazıları için sabitler tanımlanır:
\begin{lstlisting}[caption={Yeni Panel Sabitlerinin Tanımlanması}]
Public Const LBL_YEPYENI_PANEL_TITLE As String = "YEPYENİ ÖZEL ANALİZ PANELİ"
Public Const BTN_CAPTION_TO_YEPYENI  As String = "YEPYENİ ANALİZ"
\end{lstlisting}

\vspace{0.2cm}

\textbf{\textcolor{primary}{2. Adım: Panel Yapıcısı ve Çizim Prosedürü (\texttt{mod\_UI\_View})}}
\texttt{mod\_UI\_View} modülü içerisine yeni OLE Frame bileşenini ve içerisindeki kontrol elemanlarını (Label, Button, ListBox) çizen \texttt{CreateYEPYENIPanel} prosedürü yazılır:
\begin{lstlisting}[caption={CreateYEPYENIPanel Çizim Prosedürü}]
Public Sub CreateYEPYENIPanel(ws As Worksheet)
    Dim oleFrame As OLEObject, realFrame As MSForms.Frame
    Set oleFrame = ws.OLEObjects.Add(ClassType:="Forms.Frame.1", Left:=PANEL_LEFT, Top:=PANEL_TOP, Width:=PANEL_WIDTH, Height:=PANEL_HEIGHT)
    oleFrame.Name = "pan_YEPYENIPanel"
    Set realFrame = oleFrame.Object
    realFrame.Caption = ""
    realFrame.BackColor = COLOR_GRAY_BG
    
    AddControlLabel realFrame, "lbl_YepyeniTitle", LBL_YEPYENI_PANEL_TITLE, 40, 20
    AddControlListBox realFrame, "lst_YepyeniStatus", 4, "150;100;120;150"
    AddControlButton realFrame, "btn_BackToLogin", BTN_CAPTION_BACK, 40, 260, 120, 35, COLOR_RED
    AddControlButton realFrame, "btn_UpdateYepyeni", BTN_CAPTION_UPDATE, 380, 260, 120, 35, COLOR_WHITE, &H333333
    oleFrame.Visible = False
End Sub
\end{lstlisting}

\vspace{0.2cm}

\textbf{\textcolor{primary}{3. Adım: Sistem Kaydı ve Navigasyon Entegrasyonu (\texttt{mod\_UI\_Manager})}}
\begin{itemize}\setlength{\itemsep}{2pt}
    \item \texttt{BuildNavigationSystem} prosedürüne \texttt{mod\_UI\_View.CreateYEPYENIPanel ws} çağrısı eklenir.
    \item \texttt{SwitchTo} prosedüründeki panel gizleme alanına \texttt{ws.OLEObjects("pan\_YEPYENIPanel").Visible = False} satırı ve \texttt{Select Case} bloğuna \texttt{"YEPYENIPanel"} seçeneği dahil edilir.
\end{itemize}

\vspace{0.2cm}

\textbf{\textcolor{primary}{4. Adım: Olay Dinleyici ve Yönlendirme Katmanı (\texttt{clsDynamicButton} \& \texttt{mod\_UI\_Events})}}
\begin{itemize}\setlength{\itemsep}{2pt}
    \item \texttt{clsDynamicButton} sınıf modülündeki \texttt{Select Case parentFrame.Name} kontrolüne \texttt{Case "pan\_YEPYENIPanel": mod\_UI\_Events.HandleYEPYENIPanelEvents DynamicButton.Name, parentFrame} satırı eklenir.
    \item \texttt{mod\_UI\_Events} modülüne tıklamaları yöneten \texttt{HandleYEPYENIPanelEvents} prosedürü yazılır.
\end{itemize}

\newpage

% ==========================================
% SENARYO 3: YENİ HESAPLAMA SERVİSİ EKLEME
% ==========================================
\section{Senaryo 3: Yeni Bir İş Mantığı veya Hesaplama Servisi Ekleme}

Uygulamaya yeni bir hesaplama, metrik analizi veya özel rapor tarama özelliği eklenmesi gerektiğinde izlenecek mimari adımlar:

\vspace{0.3cm}

\textbf{\textcolor{primary}{1. Servis Fonksiyonunu \texttt{mod\_Services} İçinde Kurgulayın:}}
Tüm veri işleme, filtreleme ve matris kesişim adresi bulma mantığını arayüz görsel kodlarından tamamen bağımsız olarak bu modülde bir \texttt{Function} olarak yazın:
\begin{lstlisting}[caption={Yeni Hesaplama Servis Fonksiyonu Örneği}]
Public Function GetYepyeniCalculatedData() As Collection
    Dim results As New Collection, ws As Worksheet
    ' ... Veri tarama ve Scripting.Dictionary ile gruplama adimlari ...
    Set GetYepyeniCalculatedData = results
End Function
\end{lstlisting}

\vspace{0.2cm}

\textbf{\textcolor{primary}{2. Bellek İçi (In-Memory) Diziler Kullanın:}}
Excel hücrelerini döngü ile tek tek gezmek işlem hızını düşürür. Bunun yerine verileri tek hamlede \texttt{Variant} dizisine veya \texttt{Scripting.Dictionary} nesnesine çekerek RAM üzerinde işleyin:
\begin{lstlisting}[caption={RAM Üzerinde Hızlı Veri İşleme}]
Dim dataArray As Variant
dataArray = ws.Range("A2:Z1000").Value ' Tek seferde RAM'e alma
' dataArray(i, j) uzerinden milisaniyeler icinde sorgulama
\end{lstlisting}

\vspace{0.2cm}

\textbf{\textcolor{primary}{3. Sonucu Arayüz Katmanında Render Edin:}}
Hesaplanan sonuçları bir \texttt{Collection} veya dizi olarak döndürün, ardından \texttt{mod\_UI\_View} içindeki \texttt{RenderYEPYENIPanel} prosedürü ile \texttt{ListBox} elemanına bağlayın:
\begin{lstlisting}[caption={ListBox Veri Bağlama (Rendering)}]
Public Sub RenderYEPYENIPanel(realFrame As MSForms.Frame)
    Dim lst As MSForms.ListBox, data As Collection, item As Variant
    Set lst = realFrame.Controls("lst_YepyeniStatus")
    lst.Clear
    Set data = mod_Services.GetYepyeniCalculatedData()
    For Each item In data
        lst.AddItem
        lst.List(lst.ListCount - 1, 0) = item(0)
    Next item
End Sub
\end{lstlisting}

\vspace{0.3cm}

\begin{infobox}{\faSitemap\ Servis Mimarisinin Avantajı}
İş mantığı servis fonksiyonunda kurgulandığı için, gelecekte Excel arayüzü değişse bile \texttt{GetYepyeniCalculatedData} fonksiyonuna dokunulmadan farklı mecralarda tekrar kullanılabilir (Reusability).
\end{infobox}

\newpage

% ==========================================
% SENARYO 4: VERİTABANI/SAYFA EKLEME & MİMARİ KURALLAR
% ==========================================
\section{Senaryo 4: Yeni Veritabanı Çalışma Sayfası Ekleme ve Geliştirme Kuralları}

Sisteme \texttt{ARIZA} ve \texttt{BAKIM} dışında örneğin \texttt{KALITE} veya \texttt{LOJISTIK} adında 3. bir veritabanı sayfası ekleneceği zaman izlenecek süreçler ve uyulması gereken katı geliştirme kuralları şunlardır:

\vspace{0.3cm}

\textbf{\textcolor{primary}{1. Veritabanı Eşleme Fonksiyonunu Güncelleyin:}}
\texttt{mod\_Services} modülündeki \texttt{GetTargetDbSheetName} fonksiyonuna yeni sayfa koşulunu ekleyin:
\begin{lstlisting}[caption={Hedef Veritabanı Eşleme Güncellemesi}]
Public Function GetTargetDbSheetName(ByVal sheetName As String) As String
    If InStr(1, sheetName, "KALITE", vbTextCompare) > 0 Then
        GetTargetDbSheetName = "KALITE"
    ElseIf InStr(1, sheetName, "BAKIM", vbTextCompare) > 0 Then
        GetTargetDbSheetName = DATABASE_SHEET_NAME_BAKIM
    Else
        GetTargetDbSheetName = DATABASE_SHEET_NAME_ARIZA
    End If
End Function
\end{lstlisting}

\vspace{0.2cm}

\textbf{\textcolor{primary}{2. Yeni Veritabanı Sayfa Yapısını Oluşturun:}}
\begin{itemize}\setlength{\itemsep}{2pt}
    \item Sayfanın 1. satırına takip edilecek tarihleri \texttt{dd.mm.yyyy} formatında yerleştirin.
    \item B sütununa (\texttt{B} kolonu) personel/vaka isimlerini yazın.
    \item \textbf{"AÇIK VAKA SAYISI"} ve \textbf{"KAPALI VAKA SAYISI"} arama başlıklarını ilgili satırlara ekleyin.
\end{itemize}

\vspace{0.3cm}

\begin{warningbox}{\faExclamationTriangle\ Katı Kodlama Standartları ve Kalite Kuralları}
Geliştiricilerin kod tabanına müdahale ederken uymakla yükümlü olduğu temel kurallar:
\begin{itemize}\setlength{\itemsep}{2pt}
    \item \textbf{Değişken Bildirimi:} Eklenen her yeni modülün ve sınıfın en başında mutlaka \texttt{Option Explicit} ifadesi yer almalıdır.
    \item \textbf{Hücre Bazlı Döngü Yasağı:} Çalışma sayfası hücreleri üzerinde tek tek \texttt{For Each cell In Range} döngüsü kurulmamalı; veriler belleğe (\texttt{Variant Array} / \texttt{Scripting.Dictionary}) alınarak milisaniyeler seviyesinde işlenmelidir.
    \item \textbf{Bellek Güvenliği:} Dinamik oluşturulan her bir \texttt{MSForms.CommandButton} nesnesi, bellek kayıplarını ve tıklama hatalarını önlemek için merkezi \texttt{ButtonCollection} koleksiyonuna eklenmelidir.
\end{itemize}
\end{warningbox}

\vspace{0.3cm}

\begin{tcolorbox}[colback=yellow!15,colframe=orange!80!black,leftrule=3.5pt,boxrule=0.4pt,top=5pt,bottom=5pt]
    \textbf{\textcolor{orange!80!black}{\faInfoCircle\ Özet ve Katkı Yaklaşımı:}}\\
    \small Bu proje, Modüler Katmanlı Mimari (Layered Architecture) ve MVC (Model-View-Controller) prensipleri doğrultusunda inşa edilmiştir. Yapılacak tüm geliştirmelerde iş mantığı (\texttt{mod\_Services}), görünüm (\texttt{mod\_UI\_View}) ve olay yönlendirme (\texttt{mod\_UI\_Events}) katmanlarının kesinlikle birbirinden ayrık tutulması projenin sürdürülebilirliği için şarttır.
\end{tcolorbox}